% Options for packages loaded elsewhere
\PassOptionsToPackage{unicode}{hyperref}
\PassOptionsToPackage{hyphens}{url}
\documentclass[
  american,
  11pt,
  letterpaper,
]{article}
\usepackage{xcolor}
\usepackage{amsmath,amssymb}
\setcounter{secnumdepth}{5}
\usepackage{iftex}
\ifPDFTeX
  \usepackage[T1]{fontenc}
  \usepackage[utf8]{inputenc}
  \usepackage{textcomp} % provide euro and other symbols
\else % if luatex or xetex
  \usepackage{unicode-math} % this also loads fontspec
  \defaultfontfeatures{Scale=MatchLowercase}
  \defaultfontfeatures[\rmfamily]{Ligatures=TeX,Scale=1}
\fi
\usepackage{lmodern}
\ifPDFTeX\else
  % xetex/luatex font selection
\fi
% Use upquote if available, for straight quotes in verbatim environments
\IfFileExists{upquote.sty}{\usepackage{upquote}}{}
\IfFileExists{microtype.sty}{% use microtype if available
  \usepackage[]{microtype}
  \UseMicrotypeSet[protrusion]{basicmath} % disable protrusion for tt fonts
}{}
\makeatletter
\@ifundefined{KOMAClassName}{% if non-KOMA class
  \IfFileExists{parskip.sty}{%
    \usepackage{parskip}
  }{% else
    \setlength{\parindent}{0pt}
    \setlength{\parskip}{6pt plus 2pt minus 1pt}}
}{% if KOMA class
  \KOMAoptions{parskip=half}}
\makeatother
\ifLuaTeX
\usepackage[bidi=basic,shorthands=off]{babel}
\else
\usepackage[bidi=default,shorthands=off]{babel}
\fi
\ifLuaTeX
  \usepackage{selnolig} % disable illegal ligatures
\fi
\setlength{\emergencystretch}{3em} % prevent overfull lines
\providecommand{\tightlist}{%
  \setlength{\itemsep}{0pt}\setlength{\parskip}{0pt}}
\usepackage[margin=1in]{geometry}
\usepackage{setspace}
\usepackage{fancyhdr}
\usepackage{titlesec}
\usepackage[hang,flushmargin,bottom]{footmisc}
\usepackage{booktabs}
\usepackage{longtable}
\usepackage{array}
\usepackage{enumitem}
\usepackage{caption}

\setstretch{1.08}
\setlength{\parindent}{0pt}
\setlength{\parskip}{0.7em}
\setlength{\footnotesep}{0.7em}
\interfootnotelinepenalty=10000
\emergencystretch=2em
\raggedbottom

\setlist[itemize]{leftmargin=1.5em, itemsep=0.45em, topsep=0.45em}
\setlist[enumerate]{leftmargin=1.8em, itemsep=0.45em, topsep=0.45em}

\pagestyle{fancy}
\fancyhf{}
\fancyhead[L]{\nouppercase{\leftmark}}
\fancyhead[R]{AI Empire Research Papers}
\fancyfoot[C]{\thepage}
\renewcommand{\headrulewidth}{0.3pt}

\setcounter{secnumdepth}{3}
\setcounter{tocdepth}{2}
\renewcommand{\contentsname}{Paper Roadmap}

\titleformat{\section}{\Large\bfseries}{\thesection.}{0.65em}{}
\titleformat{\subsection}{\large\bfseries}{\thesubsection.}{0.6em}{}
\titleformat{\subsubsection}{\normalsize\bfseries}{\thesubsubsection.}{0.55em}{}

\captionsetup{font=small, labelfont=bf}
\AtBeginEnvironment{longtable}{\small}

\AtBeginDocument{
  \hypersetup{
    colorlinks=true,
    linkcolor=black,
    citecolor=black,
    urlcolor=blue
  }
}
\usepackage{bookmark}
\IfFileExists{xurl.sty}{\usepackage{xurl}}{} % add URL line breaks if available
\urlstyle{same}
% fallback for those not using the hyperref driver hyperxmp:
\makeatletter
\@ifundefined{xmpquote}{\newcommand{\xmpquote}[1]{#1}}{}
\makeatother
\hypersetup{
  pdftitle={Chapter 02: Social Erosion and Moral Reconfiguration},
  pdflang={en-US},
  hidelinks,
  pdfcreator={LaTeX via pandoc}}

\title{Chapter 02: Social Erosion and Moral Reconfiguration}
\usepackage{etoolbox}
\makeatletter
\providecommand{\subtitle}[1]{% add subtitle to \maketitle
  \apptocmd{\@title}{\par {\large #1 \par}}{}{}
}
\makeatother
\subtitle{From Surveillance Capitalism to AI Imperialism \textbar{}
Chapter 02: Social erosion and moral reconfiguration}
\author{}
\date{}

\begin{document}
\maketitle

\subsection{Abstract}\label{abstract}

The platform system did not only capture behavior. The U.S. If attention
is trained by systems optimized for engagement rather than flourishing,
then the outcome is not only distraction.

\subsection{Keywords}\label{keywords}

Social erosion and moral reconfiguration; surveillance capitalism; AI
imperialism; Section 230 reform; evidence-based dossier

\subsection{Research Question}\label{research-question}

How does social erosion and moral reconfiguration function within the
dossier's larger argument, and what does the documentary record allow
this chapter to establish?

\subsection{Scope and Framing Note}\label{scope-and-framing-note}

This paper examines social erosion and moral reconfiguration within the
broader dossier \emph{From Surveillance Capitalism to AI Imperialism}.
It is intentionally written to circulate on its own: necessary context
is restated, evidence boundaries are made explicit, and cited material
is reproduced in paper-local form rather than delegated to repo
navigation. It contributes to the series as a paper that deepens the
Section 230 case by connecting amplification systems to youth harm,
civic erosion, and moral reconfiguration. Adjacent reinforcement appears
in Chapter 01, Chapter 03, Chapter 10, but those cross-references are
supportive rather than required for basic comprehension.

\subsection{Central Thesis}\label{central-thesis}

The platform system did not only capture behavior. It also reorganized
norms. The strongest version of the claim is not nostalgic moral panic.
It is that algorithmic optimization degraded attention, weakened trust,
amplified social fragmentation, and altered what appears normal,
desirable, and glamorous at scale. \footnote{Sources: HHS,
  \href{https://www.hhs.gov/surgeongeneral/priorities/youth-mental-health/social-media/index.html}{Social
  Media and Youth Mental Health:\ldots{}} (n.d.; see
  \hyperref[appendix-a1-social-media-and-youth-mental-health-the-u-s-surgeon-general-s-advisory]{Appendix
  A1}); arXiv, \href{https://arxiv.org/abs/2302.03671}{What Do We Mean
  When We Talk about Tr\ldots{}} (2023-02-07; see
  \hyperref[appendix-a2-what-do-we-mean-when-we-talk-about-trust-in-social-media-a-systematic-review]{Appendix
  A2}); arXiv, \href{https://arxiv.org/abs/2010.00590}{Quantifying
  social organization and p\ldots{}} (2020-10-01; see
  \hyperref[appendix-a3-quantifying-social-organization-and-political-polarization-in-online-platforms]{Appendix
  A3})}

\subsection{Evidence and Method Note}\label{evidence-and-method-note}

This paper is derived from the chapter corpus for Chapter 02 and is
grounded in the project's structured research OS. It relies on source
notes, extracted claims, entity profiles, and timeline events already
normalized under the dossier's evidence hierarchy. Documented facts are
asserted directly where the record supports them; disputed matters are
marked as such; interpretive claims are bounded explicitly; and
speculative overreach is isolated in the evidence-boundary appendix.

\subsection{Introduction}\label{introduction}

This paper asks: How does social erosion and moral reconfiguration
function within the dossier's larger argument, and what does the
documentary record allow this chapter to establish? Its working answer
is clear from the start: The platform system did not only capture
behavior. It also reorganized norms. The strongest version of the claim
is not nostalgic moral panic. It is that algorithmic optimization
degraded attention, weakened trust, amplified social fragmentation, and
altered what appears normal, desirable, and glamorous at scale.
\footnote{Sources: HHS,
  \href{https://www.hhs.gov/surgeongeneral/priorities/youth-mental-health/social-media/index.html}{Social
  Media and Youth Mental Health:\ldots{}} (n.d.; see
  \hyperref[appendix-a1-social-media-and-youth-mental-health-the-u-s-surgeon-general-s-advisory]{Appendix
  A1}); arXiv, \href{https://arxiv.org/abs/2302.03671}{What Do We Mean
  When We Talk about Tr\ldots{}} (2023-02-07; see
  \hyperref[appendix-a2-what-do-we-mean-when-we-talk-about-trust-in-social-media-a-systematic-review]{Appendix
  A2}); arXiv, \href{https://arxiv.org/abs/2010.00590}{Quantifying
  social organization and p\ldots{}} (2020-10-01; see
  \hyperref[appendix-a3-quantifying-social-organization-and-political-polarization-in-online-platforms]{Appendix
  A3})} The objective is not only to recount developments, but to show
why the subject of this chapter belongs inside a structural argument
about extraction, enclosure, legitimacy, and reform.

\subsection{Historical or Institutional
Context}\label{historical-or-institutional-context}

The U.S. Surgeon General's advisory provides serious institutional
backing for discussing cumulative psychosocial risk among young people.
That means the youth-harm argument does not rest only on intuition or
generational panic. \footnote{Sources: HHS,
  \href{https://www.hhs.gov/surgeongeneral/priorities/youth-mental-health/social-media/index.html}{Social
  Media and Youth Mental Health:\ldots{}} (n.d.; see
  \hyperref[appendix-a1-social-media-and-youth-mental-health-the-u-s-surgeon-general-s-advisory]{Appendix
  A1})}

\subsection{Main Analysis}\label{main-analysis}

The Instagram body-image and TikTok under-13 records show that harms are
not confined to public feeds. Private-message dynamics, recommendation
mismatches, and child-directed product failures all matter. Platform
design is therefore part of the harm structure, not just the venue in
which the harm appears. \footnote{Sources: arXiv,
  \href{https://arxiv.org/abs/2504.02176}{Unfiltered: How Teens Engage
  in Body\ldots{}} (2025-04-02; see
  \hyperref[appendix-a4-unfiltered-how-teens-engage-in-body-image-and-shaming-discussions-via-instagram-direct-messages-dms]{Appendix
  A4}); arXiv, \href{https://arxiv.org/abs/2507.00299}{When Kids Mode
  Isn't For Kids: Invest\ldots{}} (2025-06-30; see
  \hyperref[appendix-a5-when-kids-mode-isn-t-for-kids-investigating-tiktok-s-under-13-experience]{Appendix
  A5})}

The trust and polarization literature strengthens the civic layer.
Social media trust, content trust, and user trust are analytically
distinct, but all are reshaped by platform architecture. Cross-national
Pew reporting adds that large majorities across surveyed countries
associate social media with manipulation, division, and reduced
civility, while the United States appears as an unusually negative
outlier. \footnote{Sources: arXiv,
  \href{https://arxiv.org/abs/2302.03671}{What Do We Mean When We Talk
  about Tr\ldots{}} (2023-02-07; see
  \hyperref[appendix-a2-what-do-we-mean-when-we-talk-about-trust-in-social-media-a-systematic-review]{Appendix
  A2}); arXiv, \href{https://arxiv.org/abs/2010.00590}{Quantifying
  social organization and p\ldots{}} (2020-10-01; see
  \hyperref[appendix-a3-quantifying-social-organization-and-political-polarization-in-online-platforms]{Appendix
  A3}); Pew Research Center,
  \href{https://www.pewresearch.org/global/2022/12/06/social-media-seen-as-mostly-good-for-democracy-across-many-nations-but-u-s-is-a-major-outlier/}{Social
  Media Seen as Mostly Good for\ldots{}} (2022-12-06; see
  \hyperref[appendix-a6-social-media-seen-as-mostly-good-for-democracy-across-many-nations-but-u-s-is-a-major-outlier]{Appendix
  A6})}

The dossier's more morally charged cultural claims also have bounded
support. Serious reporting documents recommendation systems intensifying
misogynistic content exposure. Research on OnlyFans documents mainstream
entry by creators with no prior sex-work history and ties that
normalization to platform visibility and pandemic-era adoption. These
examples should remain illustrative, but they are not empty rhetoric.
\footnote{Sources: The Guardian,
  \href{https://www.theguardian.com/media/2024/feb/06/social-media-algorithms-amplifying-misogynistic-content}{Social
  media algorithms 'amplifying m\ldots{}} (2024-02-06; see
  \hyperref[appendix-a7-social-media-algorithms-amplifying-misogynistic-content]{Appendix
  A7}); arXiv, \href{https://arxiv.org/abs/2205.10425}{Nudes? Shouldn't
  I charge for these?:\ldots{}} (2022-05-20; see
  \hyperref[appendix-a8-nudes-shouldn-t-i-charge-for-these-motivations-of-new-sexual-content-creators-on-onlyfans]{Appendix
  A8})}

Interpretation: Chapter \texttt{02} deepens the case for
\texttt{Refactor\ Section\ 230} because the social harms here are not
reducible to isolated user speech. They emerge from recommendation
systems, exposure design, monetized aspiration, and systemic
amplification. A sovereign-scale legal framework should therefore focus
on design responsibility and measurable harm, not only on moderation
after the fact. \footnote{Sources: HHS,
  \href{https://www.hhs.gov/surgeongeneral/priorities/youth-mental-health/social-media/index.html}{Social
  Media and Youth Mental Health:\ldots{}} (n.d.; see
  \hyperref[appendix-a1-social-media-and-youth-mental-health-the-u-s-surgeon-general-s-advisory]{Appendix
  A1}); arXiv, \href{https://arxiv.org/abs/2507.00299}{When Kids Mode
  Isn't For Kids: Invest\ldots{}} (2025-06-30; see
  \hyperref[appendix-a5-when-kids-mode-isn-t-for-kids-investigating-tiktok-s-under-13-experience]{Appendix
  A5}); Pew Research Center,
  \href{https://www.pewresearch.org/global/2022/12/06/social-media-seen-as-mostly-good-for-democracy-across-many-nations-but-u-s-is-a-major-outlier/}{Social
  Media Seen as Mostly Good for\ldots{}} (2022-12-06; see
  \hyperref[appendix-a6-social-media-seen-as-mostly-good-for-democracy-across-many-nations-but-u-s-is-a-major-outlier]{Appendix
  A6})}

\subsection{Position Within the
Series}\label{position-within-the-series}

This paper sits inside a coordinated 11-paper architecture. Its nearest
companions are Chapter 01 (Social networks as infrastructure for free
data capture), Chapter 03 (COVID as a historical accelerator), Chapter
10 (What to do). In that series logic, Chapter 02 deepens the Section
230 case by connecting amplification systems to youth harm, civic
erosion, and moral reconfiguration.

\subsection{Counterarguments}\label{counterarguments}

This chapter should not imply that every moral deformation of
contemporary life is fully caused by platforms. The evidence is
strongest on youth harm, fragmentation, misogynistic amplification, and
monetized sexual self-presentation. Broader civilizational claims remain
interpretive and should stay visibly bounded. \footnote{Sources: HHS,
  \href{https://www.hhs.gov/surgeongeneral/priorities/youth-mental-health/social-media/index.html}{Social
  Media and Youth Mental Health:\ldots{}} (n.d.; see
  \hyperref[appendix-a1-social-media-and-youth-mental-health-the-u-s-surgeon-general-s-advisory]{Appendix
  A1}); The Guardian,
  \href{https://www.theguardian.com/media/2024/feb/06/social-media-algorithms-amplifying-misogynistic-content}{Social
  media algorithms 'amplifying m\ldots{}} (2024-02-06; see
  \hyperref[appendix-a7-social-media-algorithms-amplifying-misogynistic-content]{Appendix
  A7})}

\subsection{Limits of the Evidence}\label{limits-of-the-evidence}

This chapter should not imply that every moral deformation of
contemporary life is fully caused by platforms. The evidence is
strongest on youth harm, fragmentation, misogynistic amplification, and
monetized sexual self-presentation. Broader civilizational claims remain
interpretive and should stay visibly bounded. \footnote{Sources: HHS,
  \href{https://www.hhs.gov/surgeongeneral/priorities/youth-mental-health/social-media/index.html}{Social
  Media and Youth Mental Health:\ldots{}} (n.d.; see
  \hyperref[appendix-a1-social-media-and-youth-mental-health-the-u-s-surgeon-general-s-advisory]{Appendix
  A1}); The Guardian,
  \href{https://www.theguardian.com/media/2024/feb/06/social-media-algorithms-amplifying-misogynistic-content}{Social
  media algorithms 'amplifying m\ldots{}} (2024-02-06; see
  \hyperref[appendix-a7-social-media-algorithms-amplifying-misogynistic-content]{Appendix
  A7})}

\subsection{Reform Relevance}\label{reform-relevance}

If attention is trained by systems optimized for engagement rather than
flourishing, then the outcome is not only distraction. It is a weaker
public with thinner trust, more manipulable norms, and reduced civic
resilience. That matters because a society softened in this way is
easier to govern through later layers of digital mediation and AI
dependency. \footnote{Sources: arXiv,
  \href{https://arxiv.org/abs/2302.03671}{What Do We Mean When We Talk
  about Tr\ldots{}} (2023-02-07; see
  \hyperref[appendix-a2-what-do-we-mean-when-we-talk-about-trust-in-social-media-a-systematic-review]{Appendix
  A2}); arXiv, \href{https://arxiv.org/abs/2010.00590}{Quantifying
  social organization and p\ldots{}} (2020-10-01; see
  \hyperref[appendix-a3-quantifying-social-organization-and-political-polarization-in-online-platforms]{Appendix
  A3}); Pew Research Center,
  \href{https://www.pewresearch.org/global/2022/12/06/social-media-seen-as-mostly-good-for-democracy-across-many-nations-but-u-s-is-a-major-outlier/}{Social
  Media Seen as Mostly Good for\ldots{}} (2022-12-06; see
  \hyperref[appendix-a6-social-media-seen-as-mostly-good-for-democracy-across-many-nations-but-u-s-is-a-major-outlier]{Appendix
  A6})}

\subsection{Conclusion}\label{conclusion}

If attention is trained by systems optimized for engagement rather than
flourishing, then the outcome is not only distraction. The next
historical shock did not create this architecture. In the architecture
of this series, Chapter 02 therefore functions as a self-contained
argument while also advancing the cumulative path toward the dossier's
three reforms.

\subsection{Notes}\label{notes}

This paper is designed to circulate independently. Public notes point
readers to the real external documents first, then to the paper-local
appendix entry that explains why each source matters.

The underlying research OS distinguishes among
\texttt{Documented\ Fact}, \texttt{Disputed\ Fact},
\texttt{Hypothesis\ /\ Interpretation}, and
\texttt{Speculative\ Narrative\ Risk}. In Chapter 02, those boundaries
remain visible in the prose, in the bibliography, and in the appendix
sections that summarize claims and caution points.

\subsection{References / Bibliography}\label{references-bibliography}

\begin{enumerate}
\def\labelenumi{\arabic{enumi}.}
\tightlist
\item
  HHS.
  \href{https://www.hhs.gov/surgeongeneral/priorities/youth-mental-health/social-media/index.html}{Social
  Media and Youth Mental Health: The U.S. Surgeon General's Advisory}.
  n.d.. T1 advisory. Paper appendix:
  \hyperref[appendix-a1-social-media-and-youth-mental-health-the-u-s-surgeon-general-s-advisory]{Appendix
  A1}. Corpus companion entry: Source S-069. Audit ID:
  \texttt{hhs-social-media-youth-advisory-2023}. Relevance: That there
  is serious institutional concern about the impact of social networks
  on minors.
\item
  arXiv. \href{https://arxiv.org/abs/2302.03671}{What Do We Mean When We
  Talk about Trust in Social Media? A Systematic Review}. 2023-02-07. T1
  paper. Paper appendix:
  \hyperref[appendix-a2-what-do-we-mean-when-we-talk-about-trust-in-social-media-a-systematic-review]{Appendix
  A2}. Corpus companion entry: Source S-120. Audit ID:
  \texttt{trust-social-media-systematic-review-2023}. Relevance: That
  trust in social media is a serious research problem with enough
  literature to support systematic synthesis.
\item
  arXiv. \href{https://arxiv.org/abs/2010.00590}{Quantifying social
  organization and political polarization in online platforms}.
  2020-10-01. T1 paper. Paper appendix:
  \hyperref[appendix-a3-quantifying-social-organization-and-political-polarization-in-online-platforms]{Appendix
  A3}. Corpus companion entry: Source S-128. Audit ID:
  \texttt{waller-anderson-online-polarization-2020}. Relevance: That
  large-scale platform social organization can be measured as fragmented
  and politically polarized rather than discussed only
  impressionistically.
\item
  arXiv. \href{https://arxiv.org/abs/2504.02176}{Unfiltered: How Teens
  Engage in Body Image and Shaming Discussions via Instagram Direct
  Messages (DMs)}. 2025-04-02. T1 paper. Paper appendix:
  \hyperref[appendix-a4-unfiltered-how-teens-engage-in-body-image-and-shaming-discussions-via-instagram-direct-messages-dms]{Appendix
  A4}. Corpus companion entry: Source S-072. Audit ID:
  \texttt{instagram-body-image-dms-2025}. Relevance: That teen
  body-image discourse on Instagram is not confined to public feeds; it
  also unfolds in private channels that shape social reinforcement and
  harm.
\item
  arXiv. \href{https://arxiv.org/abs/2507.00299}{When Kids Mode Isn't
  For Kids: Investigating TikTok's Under 13 Experience}. 2025-06-30. T1
  paper. Paper appendix:
  \hyperref[appendix-a5-when-kids-mode-isn-t-for-kids-investigating-tiktok-s-under-13-experience]{Appendix
  A5}. Corpus companion entry: Source S-118. Audit ID:
  \texttt{tiktok-kids-mode-2025}. Relevance: That a child-directed
  platform mode can still expose minors to content that the researchers
  did not classify as child-directed.
\item
  Pew Research Center.
  \href{https://www.pewresearch.org/global/2022/12/06/social-media-seen-as-mostly-good-for-democracy-across-many-nations-but-u-s-is-a-major-outlier/}{Social
  Media Seen as Mostly Good for Democracy Across Many Nations, But U.S.
  is a Major Outlier}. 2022-12-06. T2 report. Paper appendix:
  \hyperref[appendix-a6-social-media-seen-as-mostly-good-for-democracy-across-many-nations-but-u-s-is-a-major-outlier]{Appendix
  A6}. Corpus companion entry: Source S-107. Audit ID:
  \texttt{pew-social-media-democracy-outlier-2022}. Relevance: That
  cross-national evidence supports the claim that social media is widely
  associated with manipulation, division, and declining civility, even
  where some respondents still see\ldots{}
\item
  The Guardian.
  \href{https://www.theguardian.com/media/2024/feb/06/social-media-algorithms-amplifying-misogynistic-content}{Social
  media algorithms `amplifying misogynistic content'}. 2024-02-06. T2
  news\_article. Paper appendix:
  \hyperref[appendix-a7-social-media-algorithms-amplifying-misogynistic-content]{Appendix
  A7}. Corpus companion entry: Source S-067. Audit ID:
  \texttt{guardian-tiktok-misogyny-amplification-2024}. Relevance: That
  serious reporting documents a research finding that recommendation
  systems can intensify misogynistic content exposure over a short
  period.
\item
  arXiv. \href{https://arxiv.org/abs/2205.10425}{Nudes? Shouldn't I
  charge for these?: Motivations of New Sexual Content Creators on
  OnlyFans}. 2022-05-20. T1 paper. Paper appendix:
  \hyperref[appendix-a8-nudes-shouldn-t-i-charge-for-these-motivations-of-new-sexual-content-creators-on-onlyfans]{Appendix
  A8}. Corpus companion entry: Source S-068. Audit ID:
  \texttt{hamilton-onlyfans-new-creators-2022}. Relevance: That
  mainstream visibility and platform design helped normalize monetized
  sexual content creation among people with no prior sex-work history.
\end{enumerate}

\subsection{Appendix A. Source Register for This
Paper}\label{appendix-a.-source-register-for-this-paper}

\subsubsection{Appendix A1. Social Media and Youth Mental Health: The
U.S. Surgeon General's
Advisory}\label{appendix-a1-social-media-and-youth-mental-health-the-u-s-surgeon-general-s-advisory}

\textbf{Source citation:}
\href{https://www.hhs.gov/surgeongeneral/priorities/youth-mental-health/social-media/index.html}{Social
Media and Youth Mental Health: The U.S. Surgeon General's Advisory}
\textbf{Institution / author:} HHS \textbf{Date:} n.d. \textbf{Tier /
type:} T1 / advisory \textbf{Why it matters in this paper:} That there
is serious institutional concern about the impact of social networks on
minors. \textbf{Corpus companion entry:} Source S-069 \textbf{Audit ID:}
\texttt{hhs-social-media-youth-advisory-2023}

\subsubsection{Appendix A2. What Do We Mean When We Talk about Trust in
Social Media? A Systematic
Review}\label{appendix-a2-what-do-we-mean-when-we-talk-about-trust-in-social-media-a-systematic-review}

\textbf{Source citation:} \href{https://arxiv.org/abs/2302.03671}{What
Do We Mean When We Talk about Trust in Social Media? A Systematic
Review} \textbf{Institution / author:} arXiv \textbf{Date:} 2023-02-07
\textbf{Tier / type:} T1 / paper \textbf{Why it matters in this paper:}
That trust in social media is a serious research problem with enough
literature to support systematic synthesis. \textbf{Corpus companion
entry:} Source S-120 \textbf{Audit ID:}
\texttt{trust-social-media-systematic-review-2023}

\subsubsection{Appendix A3. Quantifying social organization and
political polarization in online
platforms}\label{appendix-a3-quantifying-social-organization-and-political-polarization-in-online-platforms}

\textbf{Source citation:}
\href{https://arxiv.org/abs/2010.00590}{Quantifying social organization
and political polarization in online platforms} \textbf{Institution /
author:} arXiv \textbf{Date:} 2020-10-01 \textbf{Tier / type:} T1 /
paper \textbf{Why it matters in this paper:} That large-scale platform
social organization can be measured as fragmented and politically
polarized rather than discussed only impressionistically. \textbf{Corpus
companion entry:} Source S-128 \textbf{Audit ID:}
\texttt{waller-anderson-online-polarization-2020}

\subsubsection{Appendix A4. Unfiltered: How Teens Engage in Body Image
and Shaming Discussions via Instagram Direct Messages
(DMs)}\label{appendix-a4-unfiltered-how-teens-engage-in-body-image-and-shaming-discussions-via-instagram-direct-messages-dms}

\textbf{Source citation:}
\href{https://arxiv.org/abs/2504.02176}{Unfiltered: How Teens Engage in
Body Image and Shaming Discussions via Instagram Direct Messages (DMs)}
\textbf{Institution / author:} arXiv \textbf{Date:} 2025-04-02
\textbf{Tier / type:} T1 / paper \textbf{Why it matters in this paper:}
That teen body-image discourse on Instagram is not confined to public
feeds; it also unfolds in private channels that shape social
reinforcement and harm. \textbf{Corpus companion entry:} Source S-072
\textbf{Audit ID:} \texttt{instagram-body-image-dms-2025}

\subsubsection{Appendix A5. When Kids Mode Isn't For Kids: Investigating
TikTok's Under 13
Experience}\label{appendix-a5-when-kids-mode-isn-t-for-kids-investigating-tiktok-s-under-13-experience}

\textbf{Source citation:} \href{https://arxiv.org/abs/2507.00299}{When
Kids Mode Isn't For Kids: Investigating TikTok's Under 13 Experience}
\textbf{Institution / author:} arXiv \textbf{Date:} 2025-06-30
\textbf{Tier / type:} T1 / paper \textbf{Why it matters in this paper:}
That a child-directed platform mode can still expose minors to content
that the researchers did not classify as child-directed. \textbf{Corpus
companion entry:} Source S-118 \textbf{Audit ID:}
\texttt{tiktok-kids-mode-2025}

\subsubsection{Appendix A6. Social Media Seen as Mostly Good for
Democracy Across Many Nations, But U.S. is a Major
Outlier}\label{appendix-a6-social-media-seen-as-mostly-good-for-democracy-across-many-nations-but-u-s-is-a-major-outlier}

\textbf{Source citation:}
\href{https://www.pewresearch.org/global/2022/12/06/social-media-seen-as-mostly-good-for-democracy-across-many-nations-but-u-s-is-a-major-outlier/}{Social
Media Seen as Mostly Good for Democracy Across Many Nations, But U.S. is
a Major Outlier} \textbf{Institution / author:} Pew Research Center
\textbf{Date:} 2022-12-06 \textbf{Tier / type:} T2 / report \textbf{Why
it matters in this paper:} That cross-national evidence supports the
claim that social media is widely associated with manipulation,
division, and declining civility, even where some respondents still
see\ldots{} \textbf{Corpus companion entry:} Source S-107 \textbf{Audit
ID:} \texttt{pew-social-media-democracy-outlier-2022}

\subsubsection{Appendix A7. Social media algorithms `amplifying
misogynistic
content'}\label{appendix-a7-social-media-algorithms-amplifying-misogynistic-content}

\textbf{Source citation:}
\href{https://www.theguardian.com/media/2024/feb/06/social-media-algorithms-amplifying-misogynistic-content}{Social
media algorithms `amplifying misogynistic content'} \textbf{Institution
/ author:} The Guardian \textbf{Date:} 2024-02-06 \textbf{Tier / type:}
T2 / news\_article \textbf{Why it matters in this paper:} That serious
reporting documents a research finding that recommendation systems can
intensify misogynistic content exposure over a short period.
\textbf{Corpus companion entry:} Source S-067 \textbf{Audit ID:}
\texttt{guardian-tiktok-misogyny-amplification-2024}

\subsubsection{Appendix A8. Nudes? Shouldn't I charge for these?:
Motivations of New Sexual Content Creators on
OnlyFans}\label{appendix-a8-nudes-shouldn-t-i-charge-for-these-motivations-of-new-sexual-content-creators-on-onlyfans}

\textbf{Source citation:} \href{https://arxiv.org/abs/2205.10425}{Nudes?
Shouldn't I charge for these?: Motivations of New Sexual Content
Creators on OnlyFans} \textbf{Institution / author:} arXiv
\textbf{Date:} 2022-05-20 \textbf{Tier / type:} T1 / paper \textbf{Why
it matters in this paper:} That mainstream visibility and platform
design helped normalize monetized sexual content creation among people
with no prior sex-work history. \textbf{Corpus companion entry:} Source
S-068 \textbf{Audit ID:} \texttt{hamilton-onlyfans-new-creators-2022}

\subsection{Appendix B. Claims Used in This
Paper}\label{appendix-b.-claims-used-in-this-paper}

\subsubsection{Appendix B1. Serious reporting documents a research
finding that recommendation systems can
in\ldots{}}\label{appendix-b1-serious-reporting-documents-a-research-finding-that-recommendation-systems-can-in}

\textbf{Claim statement:} That serious reporting documents a research
finding that recommendation systems can intensify misogynistic content
exposure over a short period. \textbf{Evidence status:} Documented Fact
\textbf{Source support:}
\hyperref[appendix-a7-social-media-algorithms-amplifying-misogynistic-content]{Appendix
A7} \textbf{Corpus companion entry:} Claim C-210 \textbf{Audit Claim
ID:}
\texttt{guardian-tiktok-misogyny-amplification-2024-02-social-decay-youth-harm-hypersexualization-claim-01}

\subsubsection{Appendix B2. Chapter 2 can discuss algorithmic
amplification of misogynistic norm formation
wi\ldots{}}\label{appendix-b2-chapter-2-can-discuss-algorithmic-amplification-of-misogynistic-norm-formation-wi}

\textbf{Claim statement:} That Chapter 2 can discuss algorithmic
amplification of misogynistic norm formation with more than anecdotal
support. \textbf{Evidence status:} Documented Fact \textbf{Source
support:}
\hyperref[appendix-a7-social-media-algorithms-amplifying-misogynistic-content]{Appendix
A7} \textbf{Corpus companion entry:} Claim C-211 \textbf{Audit Claim
ID:}
\texttt{guardian-tiktok-misogyny-amplification-2024-02-social-decay-youth-harm-hypersexualization-claim-02}

\subsubsection{Appendix B3. The dossier can connect youth-facing
platform design to moral and social
reconfig\ldots{}}\label{appendix-b3-the-dossier-can-connect-youth-facing-platform-design-to-moral-and-social-reconfig}

\textbf{Claim statement:} That the dossier can connect youth-facing
platform design to moral and social reconfiguration without reducing the
argument to a generalized complaint about bad culture. \textbf{Evidence
status:} Documented Fact \textbf{Source support:}
\hyperref[appendix-a7-social-media-algorithms-amplifying-misogynistic-content]{Appendix
A7} \textbf{Corpus companion entry:} Claim C-212 \textbf{Audit Claim
ID:}
\texttt{guardian-tiktok-misogyny-amplification-2024-02-social-decay-youth-harm-hypersexualization-claim-03}

\subsubsection{Appendix B4. Mainstream visibility and platform design
helped normalize monetized sexual
conte\ldots{}}\label{appendix-b4-mainstream-visibility-and-platform-design-helped-normalize-monetized-sexual-conte}

\textbf{Claim statement:} That mainstream visibility and platform design
helped normalize monetized sexual content creation among people with no
prior sex-work history. \textbf{Evidence status:} Documented Fact
\textbf{Source support:}
\hyperref[appendix-a8-nudes-shouldn-t-i-charge-for-these-motivations-of-new-sexual-content-creators-on-onlyfans]{Appendix
A8} \textbf{Corpus companion entry:} Claim C-213 \textbf{Audit Claim
ID:}
\texttt{hamilton-onlyfans-new-creators-2022-02-social-decay-youth-harm-hypersexualization-claim-01}

\subsubsection{Appendix B5. The pandemic was a documented adoption
driver in this
transition}\label{appendix-b5-the-pandemic-was-a-documented-adoption-driver-in-this-transition}

\textbf{Claim statement:} That the pandemic was a documented adoption
driver in this transition. \textbf{Evidence status:} Documented Fact
\textbf{Source support:}
\hyperref[appendix-a8-nudes-shouldn-t-i-charge-for-these-motivations-of-new-sexual-content-creators-on-onlyfans]{Appendix
A8} \textbf{Corpus companion entry:} Claim C-214 \textbf{Audit Claim
ID:}
\texttt{hamilton-onlyfans-new-creators-2022-02-social-decay-youth-harm-hypersexualization-claim-02}

\subsubsection{Appendix B6. Chapter 2 can discuss digital prostitution
glamour and normalization with a
real\ldots{}}\label{appendix-b6-chapter-2-can-discuss-digital-prostitution-glamour-and-normalization-with-a-real}

\textbf{Claim statement:} That Chapter 2 can discuss digital
prostitution glamour and normalization with a real research basis rather
than relying only on celebrity examples or moral impression.
\textbf{Evidence status:} Documented Fact \textbf{Source support:}
\hyperref[appendix-a8-nudes-shouldn-t-i-charge-for-these-motivations-of-new-sexual-content-creators-on-onlyfans]{Appendix
A8} \textbf{Corpus companion entry:} Claim C-215 \textbf{Audit Claim
ID:}
\texttt{hamilton-onlyfans-new-creators-2022-02-social-decay-youth-harm-hypersexualization-claim-03}

\subsubsection{Appendix B7. There is serious institutional concern about
the impact of social networks on
minors}\label{appendix-b7-there-is-serious-institutional-concern-about-the-impact-of-social-networks-on-minors}

\textbf{Claim statement:} That there is serious institutional concern
about the impact of social networks on minors. \textbf{Evidence status:}
Documented Fact \textbf{Source support:}
\hyperref[appendix-a1-social-media-and-youth-mental-health-the-u-s-surgeon-general-s-advisory]{Appendix
A1} \textbf{Corpus companion entry:} Claim C-216 \textbf{Audit Claim
ID:} \texttt{hhs-social-media-youth-advisory-2023-claim-01}

\subsubsection{Appendix B8. There is enough basis to discuss
psychosocial harm and cumulative risks in
young\ldots{}}\label{appendix-b8-there-is-enough-basis-to-discuss-psychosocial-harm-and-cumulative-risks-in-young}

\textbf{Claim statement:} That there is enough basis to discuss
psychosocial harm and cumulative risks in young populations.
\textbf{Evidence status:} Documented Fact \textbf{Source support:}
\hyperref[appendix-a1-social-media-and-youth-mental-health-the-u-s-surgeon-general-s-advisory]{Appendix
A1} \textbf{Corpus companion entry:} Claim C-217 \textbf{Audit Claim
ID:} \texttt{hhs-social-media-youth-advisory-2023-claim-02}

\subsubsection{Appendix B9. The discussion about youth does not depend
only on moral intuition or cultural
no\ldots{}}\label{appendix-b9-the-discussion-about-youth-does-not-depend-only-on-moral-intuition-or-cultural-no}

\textbf{Claim statement:} That the discussion about youth does not
depend only on moral intuition or cultural nostalgia. \textbf{Evidence
status:} Documented Fact \textbf{Source support:}
\hyperref[appendix-a1-social-media-and-youth-mental-health-the-u-s-surgeon-general-s-advisory]{Appendix
A1} \textbf{Corpus companion entry:} Claim C-218 \textbf{Audit Claim
ID:} \texttt{hhs-social-media-youth-advisory-2023-claim-03}

\subsubsection{Appendix B10. Teen body-image discourse on Instagram is
not confined to public feeds; it also
u\ldots{}}\label{appendix-b10-teen-body-image-discourse-on-instagram-is-not-confined-to-public-feeds-it-also-u}

\textbf{Claim statement:} That teen body-image discourse on Instagram is
not confined to public feeds; it also unfolds in private channels that
shape social reinforcement and harm. \textbf{Evidence status:}
Documented Fact \textbf{Source support:}
\hyperref[appendix-a4-unfiltered-how-teens-engage-in-body-image-and-shaming-discussions-via-instagram-direct-messages-dms]{Appendix
A4} \textbf{Corpus companion entry:} Claim C-225 \textbf{Audit Claim
ID:}
\texttt{instagram-body-image-dms-2025-02-social-decay-youth-harm-hypersexualization-claim-01}

\subsubsection{Appendix B11. Platform architecture matters: the
difference between one-on-one and group
settin\ldots{}}\label{appendix-b11-platform-architecture-matters-the-difference-between-one-on-one-and-group-settin}

\textbf{Claim statement:} That platform architecture matters: the
difference between one-on-one and group settings changes the texture of
support, criticism, and shaming. \textbf{Evidence status:} Documented
Fact \textbf{Source support:}
\hyperref[appendix-a4-unfiltered-how-teens-engage-in-body-image-and-shaming-discussions-via-instagram-direct-messages-dms]{Appendix
A4} \textbf{Corpus companion entry:} Claim C-226 \textbf{Audit Claim
ID:}
\texttt{instagram-body-image-dms-2025-02-social-decay-youth-harm-hypersexualization-claim-02}

\subsubsection{Appendix B12. Chapter 2 can discuss mediated body-image
norms with stronger empirical
grounding\ldots{}}\label{appendix-b12-chapter-2-can-discuss-mediated-body-image-norms-with-stronger-empirical-grounding}

\textbf{Claim statement:} That Chapter 2 can discuss mediated body-image
norms with stronger empirical grounding than a general moral complaint.
\textbf{Evidence status:} Documented Fact \textbf{Source support:}
\hyperref[appendix-a4-unfiltered-how-teens-engage-in-body-image-and-shaming-discussions-via-instagram-direct-messages-dms]{Appendix
A4} \textbf{Corpus companion entry:} Claim C-227 \textbf{Audit Claim
ID:}
\texttt{instagram-body-image-dms-2025-02-social-decay-youth-harm-hypersexualization-claim-03}

\subsubsection{Appendix B13. Cross-national evidence supports the claim
that social media is widely
associated\ldots{}}\label{appendix-b13-cross-national-evidence-supports-the-claim-that-social-media-is-widely-associated}

\textbf{Claim statement:} That cross-national evidence supports the
claim that social media is widely associated with manipulation,
division, and declining civility, even where some respondents still see
democratic benefits. \textbf{Evidence status:} Documented Fact
\textbf{Source support:}
\hyperref[appendix-a6-social-media-seen-as-mostly-good-for-democracy-across-many-nations-but-u-s-is-a-major-outlier]{Appendix
A6} \textbf{Corpus companion entry:} Claim C-353 \textbf{Audit Claim
ID:}
\texttt{pew-social-media-democracy-outlier-2022-02-social-decay-youth-harm-hypersexualization-claim-01}

\subsubsection{Appendix B14. The United States can be treated as an
unusually severe case in the
trust-and-fra\ldots{}}\label{appendix-b14-the-united-states-can-be-treated-as-an-unusually-severe-case-in-the-trust-and-fra}

\textbf{Claim statement:} That the United States can be treated as an
unusually severe case in the trust-and-fragmentation story rather than
as just another generic example. \textbf{Evidence status:} Documented
Fact \textbf{Source support:}
\hyperref[appendix-a6-social-media-seen-as-mostly-good-for-democracy-across-many-nations-but-u-s-is-a-major-outlier]{Appendix
A6} \textbf{Corpus companion entry:} Claim C-354 \textbf{Audit Claim
ID:}
\texttt{pew-social-media-democracy-outlier-2022-02-social-decay-youth-harm-hypersexualization-claim-02}

\subsubsection{Appendix B15. Chapter 2 can connect platform-mediated
social erosion to democratic trust and
ci\ldots{}}\label{appendix-b15-chapter-2-can-connect-platform-mediated-social-erosion-to-democratic-trust-and-ci}

\textbf{Claim statement:} That Chapter 2 can connect platform-mediated
social erosion to democratic trust and civic coherence with stronger
evidence than youth-harm studies alone provide. \textbf{Evidence
status:} Documented Fact \textbf{Source support:}
\hyperref[appendix-a6-social-media-seen-as-mostly-good-for-democracy-across-many-nations-but-u-s-is-a-major-outlier]{Appendix
A6} \textbf{Corpus companion entry:} Claim C-355 \textbf{Audit Claim
ID:}
\texttt{pew-social-media-democracy-outlier-2022-02-social-decay-youth-harm-hypersexualization-claim-03}

\subsubsection{Appendix B16. A child-directed platform mode can still
expose minors to content that the
resear\ldots{}}\label{appendix-b16-a-child-directed-platform-mode-can-still-expose-minors-to-content-that-the-resear}

\textbf{Claim statement:} That a child-directed platform mode can still
expose minors to content that the researchers did not classify as
child-directed. \textbf{Evidence status:} Documented Fact \textbf{Source
support:}
\hyperref[appendix-a5-when-kids-mode-isn-t-for-kids-investigating-tiktok-s-under-13-experience]{Appendix
A5} \textbf{Corpus companion entry:} Claim C-385 \textbf{Audit Claim
ID:}
\texttt{tiktok-kids-mode-2025-02-social-decay-youth-harm-hypersexualization-claim-01}

\subsubsection{Appendix B17. Chapter 2 can discuss platform design
failures affecting minors with concrete
evi\ldots{}}\label{appendix-b17-chapter-2-can-discuss-platform-design-failures-affecting-minors-with-concrete-evi}

\textbf{Claim statement:} That Chapter 2 can discuss platform design
failures affecting minors with concrete evidence rather than only
generalized concern. \textbf{Evidence status:} Documented Fact
\textbf{Source support:}
\hyperref[appendix-a5-when-kids-mode-isn-t-for-kids-investigating-tiktok-s-under-13-experience]{Appendix
A5} \textbf{Corpus companion entry:} Claim C-386 \textbf{Audit Claim
ID:}
\texttt{tiktok-kids-mode-2025-02-social-decay-youth-harm-hypersexualization-claim-02}

\subsubsection{Appendix B18. Youth-risk arguments in the dossier can
include the mismatch between stated
safet\ldots{}}\label{appendix-b18-youth-risk-arguments-in-the-dossier-can-include-the-mismatch-between-stated-safet}

\textbf{Claim statement:} That youth-risk arguments in the dossier can
include the mismatch between stated safety framing and actual curation
outcomes. \textbf{Evidence status:} Documented Fact \textbf{Source
support:}
\hyperref[appendix-a5-when-kids-mode-isn-t-for-kids-investigating-tiktok-s-under-13-experience]{Appendix
A5} \textbf{Corpus companion entry:} Claim C-387 \textbf{Audit Claim
ID:}
\texttt{tiktok-kids-mode-2025-02-social-decay-youth-harm-hypersexualization-claim-03}

\subsubsection{Appendix B19. Trust in social media is a serious research
problem with enough literature to
sup\ldots{}}\label{appendix-b19-trust-in-social-media-is-a-serious-research-problem-with-enough-literature-to-sup}

\textbf{Claim statement:} That trust in social media is a serious
research problem with enough literature to support systematic synthesis.
\textbf{Evidence status:} Documented Fact \textbf{Source support:}
\hyperref[appendix-a2-what-do-we-mean-when-we-talk-about-trust-in-social-media-a-systematic-review]{Appendix
A2} \textbf{Corpus companion entry:} Claim C-391 \textbf{Audit Claim
ID:}
\texttt{trust-social-media-systematic-review-2023-02-social-decay-youth-harm-hypersexualization-claim-01}

\subsubsection{Appendix B20. Trust in platforms, users, and content
should not be treated as interchangeable
c\ldots{}}\label{appendix-b20-trust-in-platforms-users-and-content-should-not-be-treated-as-interchangeable-c}

\textbf{Claim statement:} That trust in platforms, users, and content
should not be treated as interchangeable concepts inside the dossier.
\textbf{Evidence status:} Documented Fact \textbf{Source support:}
\hyperref[appendix-a2-what-do-we-mean-when-we-talk-about-trust-in-social-media-a-systematic-review]{Appendix
A2} \textbf{Corpus companion entry:} Claim C-392 \textbf{Audit Claim
ID:}
\texttt{trust-social-media-systematic-review-2023-02-social-decay-youth-harm-hypersexualization-claim-02}

\subsubsection{Appendix B21. Chapter 2 can discuss erosion of trust and
social interpretation with scholarly
s\ldots{}}\label{appendix-b21-chapter-2-can-discuss-erosion-of-trust-and-social-interpretation-with-scholarly-s}

\textbf{Claim statement:} That Chapter 2 can discuss erosion of trust
and social interpretation with scholarly support, while still keeping
conceptual precision. \textbf{Evidence status:} Documented Fact
\textbf{Source support:}
\hyperref[appendix-a2-what-do-we-mean-when-we-talk-about-trust-in-social-media-a-systematic-review]{Appendix
A2} \textbf{Corpus companion entry:} Claim C-393 \textbf{Audit Claim
ID:}
\texttt{trust-social-media-systematic-review-2023-02-social-decay-youth-harm-hypersexualization-claim-03}

\subsubsection{Appendix B22. Large-scale platform social organization
can be measured as fragmented and
politi\ldots{}}\label{appendix-b22-large-scale-platform-social-organization-can-be-measured-as-fragmented-and-politi}

\textbf{Claim statement:} That large-scale platform social organization
can be measured as fragmented and politically polarized rather than
discussed only impressionistically. \textbf{Evidence status:} Documented
Fact \textbf{Source support:}
\hyperref[appendix-a3-quantifying-social-organization-and-political-polarization-in-online-platforms]{Appendix
A3} \textbf{Corpus companion entry:} Claim C-470 \textbf{Audit Claim
ID:}
\texttt{waller-anderson-online-polarization-2020-02-social-decay-youth-harm-hypersexualization-claim-01}

\subsubsection{Appendix B23. Chapter 2 can move beyond youth harm alone
and document broader social
fragmentat\ldots{}}\label{appendix-b23-chapter-2-can-move-beyond-youth-harm-alone-and-document-broader-social-fragmentat}

\textbf{Claim statement:} That Chapter 2 can move beyond youth harm
alone and document broader social fragmentation as a platform-structured
phenomenon. \textbf{Evidence status:} Documented Fact \textbf{Source
support:}
\hyperref[appendix-a3-quantifying-social-organization-and-political-polarization-in-online-platforms]{Appendix
A3} \textbf{Corpus companion entry:} Claim C-471 \textbf{Audit Claim
ID:}
\texttt{waller-anderson-online-polarization-2020-02-social-decay-youth-harm-hypersexualization-claim-02}

\subsubsection{Appendix B24. The strongest version of the argument
should focus on system-level polarization
d\ldots{}}\label{appendix-b24-the-strongest-version-of-the-argument-should-focus-on-system-level-polarization-d}

\textbf{Claim statement:} That the strongest version of the argument
should focus on system-level polarization dynamics, not a simplistic
claim that every individual user became equally radicalized.
\textbf{Evidence status:} Documented Fact \textbf{Source support:}
\hyperref[appendix-a3-quantifying-social-organization-and-political-polarization-in-online-platforms]{Appendix
A3} \textbf{Corpus companion entry:} Claim C-472 \textbf{Audit Claim
ID:}
\texttt{waller-anderson-online-polarization-2020-02-social-decay-youth-harm-hypersexualization-claim-03}

\subsubsection{Appendix B25. Documented youth exposure, recommendation
amplification, trust erosion, and
fragm\ldots{}}\label{appendix-b25-documented-youth-exposure-recommendation-amplification-trust-erosion-and-fragm}

\textbf{Claim statement:} That documented youth exposure, recommendation
amplification, trust erosion, and fragmentation support a
sovereign-scale duty framework under a refactored Section 230 focused on
design, exposure, and systemic harm. \textbf{Evidence status:}
Hypothesis / Interpretation \textbf{Source support:}
\hyperref[appendix-a1-social-media-and-youth-mental-health-the-u-s-surgeon-general-s-advisory]{Appendix
A1},
\hyperref[appendix-a6-social-media-seen-as-mostly-good-for-democracy-across-many-nations-but-u-s-is-a-major-outlier]{Appendix
A6} \textbf{Corpus companion entry:} Claim C-419 \textbf{Audit Claim
ID:} \texttt{vector02-design-duty-section230-claim-01}

\subsubsection{Appendix B26. Platform optimization alone fully explains
contemporary moral decline, every
yout\ldots{}}\label{appendix-b26-platform-optimization-alone-fully-explains-contemporary-moral-decline-every-yout}

\textbf{Claim statement:} That platform optimization alone fully
explains contemporary moral decline, every youth pathology, or every
downstream cultural change associated with digital life.
\textbf{Evidence status:} Speculative Narrative Risk \textbf{Source
support:}
\hyperref[appendix-a1-social-media-and-youth-mental-health-the-u-s-surgeon-general-s-advisory]{Appendix
A1} \textbf{Corpus companion entry:} Claim C-420 \textbf{Audit Claim
ID:} \texttt{vector02-total-moral-decline-overclaim-claim-01}

\subsection{Appendix C. Timeline
Slice}\label{appendix-c.-timeline-slice}

\subsubsection{Appendix C1. 2020-10-01: A large-scale study documents
system-level polarization dynamics across Reddit
commun\ldots{}}\label{appendix-c1-2020-10-01-a-large-scale-study-documents-system-level-polarization-dynamics-across-reddit-communities}

\textbf{Event summary:} A large-scale study documents system-level
polarization dynamics across Reddit communities \textbf{Event date:}
2020-10-01 \textbf{Source support:}
\hyperref[appendix-a3-quantifying-social-organization-and-political-polarization-in-online-platforms]{Appendix
A3} \textbf{Corpus companion entry:} Event T-008 \textbf{Audit Event
ID:} \texttt{event-reddit-polarization-study-2020-10-01}

\subsubsection{Appendix C2. 2022-05-20: A study documents mainstream
acceptance as a driver of new OnlyFans sexual-content
cr\ldots{}}\label{appendix-c2-2022-05-20-a-study-documents-mainstream-acceptance-as-a-driver-of-new-onlyfans-sexual-content-creators}

\textbf{Event summary:} A study documents mainstream acceptance as a
driver of new OnlyFans sexual-content creators \textbf{Event date:}
2022-05-20 \textbf{Source support:}
\hyperref[appendix-a8-nudes-shouldn-t-i-charge-for-these-motivations-of-new-sexual-content-creators-on-onlyfans]{Appendix
A8} \textbf{Corpus companion entry:} Event T-018 \textbf{Audit Event
ID:} \texttt{event-onlyfans-mainstreaming-study-2022-05-20}

\subsubsection{Appendix C3. 2022-12-06: Pew publishes cross-national
findings on social media, democracy, manipulation, and
d\ldots{}}\label{appendix-c3-2022-12-06-pew-publishes-cross-national-findings-on-social-media-democracy-manipulation-and-division}

\textbf{Event summary:} Pew publishes cross-national findings on social
media, democracy, manipulation, and division \textbf{Event date:}
2022-12-06 \textbf{Source support:}
\hyperref[appendix-a6-social-media-seen-as-mostly-good-for-democracy-across-many-nations-but-u-s-is-a-major-outlier]{Appendix
A6} \textbf{Corpus companion entry:} Event T-019 \textbf{Audit Event
ID:} \texttt{event-pew-social-democracy-outlier-2022-12-06}

\subsubsection{Appendix C4. 2024-02-06: Researchers report rapid
amplification of misogynistic recommendations on
TikTok}\label{appendix-c4-2024-02-06-researchers-report-rapid-amplification-of-misogynistic-recommendations-on-tiktok}

\textbf{Event summary:} Researchers report rapid amplification of
misogynistic recommendations on TikTok \textbf{Event date:} 2024-02-06
\textbf{Source support:}
\hyperref[appendix-a7-social-media-algorithms-amplifying-misogynistic-content]{Appendix
A7} \textbf{Corpus companion entry:} Event T-041 \textbf{Audit Event
ID:} \texttt{event-tiktok-misogyny-amplification-2024-02-06}

\subsection{Appendix D. Relevant
Entities}\label{appendix-d.-relevant-entities}

\subsubsection{Appendix D1. Instagram}\label{appendix-d1-instagram}

\textbf{Entity type:} platform \textbf{Role in this paper:} Actor
relevant to 02\_social\_decay\_youth\_harm\_hypersexualization.
\textbf{Relevant sources in this paper:}
\hyperref[appendix-a4-unfiltered-how-teens-engage-in-body-image-and-shaming-discussions-via-instagram-direct-messages-dms]{Appendix
A4} \textbf{Corpus companion entry:} Entity E-021 \textbf{Audit Entity
ID:} \texttt{org-instagram}

\subsubsection{Appendix D2. OnlyFans}\label{appendix-d2-onlyfans}

\textbf{Entity type:} platform \textbf{Role in this paper:} Actor
relevant to 02\_social\_decay\_youth\_harm\_hypersexualization.
\textbf{Relevant sources in this paper:}
\hyperref[appendix-a8-nudes-shouldn-t-i-charge-for-these-motivations-of-new-sexual-content-creators-on-onlyfans]{Appendix
A8} \textbf{Corpus companion entry:} Entity E-026 \textbf{Audit Entity
ID:} \texttt{org-onlyfans}

\subsubsection{Appendix D3. Reddit}\label{appendix-d3-reddit}

\textbf{Entity type:} platform \textbf{Role in this paper:} Actor
relevant to 02\_social\_decay\_youth\_harm\_hypersexualization,
04\_data\_to\_models\_copyright\_weights. \textbf{Relevant sources in
this paper:}
\hyperref[appendix-a3-quantifying-social-organization-and-political-polarization-in-online-platforms]{Appendix
A3} \textbf{Corpus companion entry:} Entity E-032 \textbf{Audit Entity
ID:} \texttt{org-reddit}

\subsubsection{Appendix D4. TikTok}\label{appendix-d4-tiktok}

\textbf{Entity type:} platform \textbf{Role in this paper:} Actor
relevant to 02\_social\_decay\_youth\_harm\_hypersexualization.
\textbf{Relevant sources in this paper:}
\hyperref[appendix-a5-when-kids-mode-isn-t-for-kids-investigating-tiktok-s-under-13-experience]{Appendix
A5},
\hyperref[appendix-a7-social-media-algorithms-amplifying-misogynistic-content]{Appendix
A7} \textbf{Corpus companion entry:} Entity E-034 \textbf{Audit Entity
ID:} \texttt{org-tiktok}

\subsubsection{Appendix D5. U.S. Surgeon General /
HHS}\label{appendix-d5-u-s-surgeon-general-hhs}

\textbf{Entity type:} public\_health \textbf{Role in this paper:} Actor
relevant to 02\_social\_decay\_youth\_harm\_hypersexualization.
\textbf{Relevant sources in this paper:}
\hyperref[appendix-a1-social-media-and-youth-mental-health-the-u-s-surgeon-general-s-advisory]{Appendix
A1} \textbf{Corpus companion entry:} Entity E-019 \textbf{Audit Entity
ID:} \texttt{org-hhs-surgeon-general}

\subsubsection{Appendix D6. Pew Research
Center}\label{appendix-d6-pew-research-center}

\textbf{Entity type:} research\_institution \textbf{Role in this paper:}
Actor relevant to 02\_social\_decay\_youth\_harm\_hypersexualization.
\textbf{Relevant sources in this paper:}
\hyperref[appendix-a6-social-media-seen-as-mostly-good-for-democracy-across-many-nations-but-u-s-is-a-major-outlier]{Appendix
A6} \textbf{Corpus companion entry:} Entity E-030 \textbf{Audit Entity
ID:} \texttt{org-pew-research-center}

\subsection{Appendix E. Evidence
Boundaries}\label{appendix-e.-evidence-boundaries}

\subsubsection{Appendix E1. Disputed
Matters}\label{appendix-e1-disputed-matters}

\begin{itemize}
\tightlist
\item
  None registered for this paper.
\end{itemize}

\subsubsection{Appendix E2. Interpretive
Boundaries}\label{appendix-e2-interpretive-boundaries}

\begin{itemize}
\tightlist
\item
  \hyperref[appendix-b25-documented-youth-exposure-recommendation-amplification-trust-erosion-and-fragm]{Appendix
  B25}: That documented youth exposure, recommendation amplification,
  trust erosion, and fragmentation support a sovereign-scale duty
  framework under a refactored Section 230 focused on design, exposure,
  and systemic harm.
\end{itemize}

\subsubsection{Appendix E3. Speculative Narrative
Risks}\label{appendix-e3-speculative-narrative-risks}

\begin{itemize}
\tightlist
\item
  \hyperref[appendix-b26-platform-optimization-alone-fully-explains-contemporary-moral-decline-every-yout]{Appendix
  B26}: That platform optimization alone fully explains contemporary
  moral decline, every youth pathology, or every downstream cultural
  change associated with digital life.
\end{itemize}

\subsubsection{Appendix E4. Sensitive or Review-Worthy Note
Flags}\label{appendix-e4-sensitive-or-review-worthy-note-flags}

\begin{itemize}
\tightlist
\item
  \hyperref[appendix-a7-social-media-algorithms-amplifying-misogynistic-content]{Appendix
  A7}. Companion note: Note N-008. Risk flag: medium. Note focus: That
  serious reporting documents a research finding that recommendation
  systems can intensify misogynistic content exposure over a short
  period.
\item
  \hyperref[appendix-a8-nudes-shouldn-t-i-charge-for-these-motivations-of-new-sexual-content-creators-on-onlyfans]{Appendix
  A8}. Companion note: Note N-009. Risk flag: medium. Note focus: That
  mainstream visibility and platform design helped normalize monetized
  sexual content creation among people with no prior sex-work history.
\item
  \hyperref[appendix-a1-social-media-and-youth-mental-health-the-u-s-surgeon-general-s-advisory]{Appendix
  A1}. Companion note: Note N-010. Risk flag: low. Note focus: That
  there is serious institutional concern about the impact of social
  networks on minors.
\item
  \hyperref[appendix-a4-unfiltered-how-teens-engage-in-body-image-and-shaming-discussions-via-instagram-direct-messages-dms]{Appendix
  A4}. Companion note: Note N-011. Risk flag: review. Note focus: That
  teen body-image discourse on Instagram is not confined to public
  feeds; it also unfolds in private channels that shape social
  reinforcement and harm.
\item
  \hyperref[appendix-a6-social-media-seen-as-mostly-good-for-democracy-across-many-nations-but-u-s-is-a-major-outlier]{Appendix
  A6}. Companion note: Note N-012. Risk flag: review. Note focus: That
  cross-national evidence supports the claim that social media is widely
  associated with manipulation, division, and declining civility, even
  where some respondents still see democratic benefits.
\item
  \hyperref[appendix-a5-when-kids-mode-isn-t-for-kids-investigating-tiktok-s-under-13-experience]{Appendix
  A5}. Companion note: Note N-013. Risk flag: review. Note focus: That a
  child-directed platform mode can still expose minors to content that
  the researchers did not classify as child-directed.
\item
  \hyperref[appendix-a2-what-do-we-mean-when-we-talk-about-trust-in-social-media-a-systematic-review]{Appendix
  A2}. Companion note: Note N-014. Risk flag: review. Note focus: That
  trust in social media is a serious research problem with enough
  literature to support systematic synthesis.
\item
  \hyperref[appendix-a3-quantifying-social-organization-and-political-polarization-in-online-platforms]{Appendix
  A3}. Companion note: Note N-015. Risk flag: medium. Note focus: That
  large-scale platform social organization can be measured as fragmented
  and politically polarized rather than discussed only
  impressionistically.
\end{itemize}

\end{document}
